We now apply \Cref{lem:repair} to the rounded labels of \Cref{ker:angles}.
The following proposition collects the metric properties needed for
the contraction and duality arguments.

\begin{proposition}\label{prop:metric}
For every sufficiently large integer $m$, the vertex set $V$ in
\eqref{ker:vertices} admits a negative-type semimetric $d$ with the
following properties. Its chart maps $v_s(\theta)$, defined almost
everywhere on $U=(-2,2)^m$, satisfy
\begin{equation}\label{met:eq:macro}
\left|d(v_s(\theta),v_t(\eta))
-c_*\sqrt m\,|\theta-\eta|\right|\le Cr
\qquad(\theta,\eta\in U)
\end{equation}
whenever the labels are defined, where $r=m^{-8}$.
For labels in the same chart,
\begin{equation}\label{met:eq:adjacent}
\|x-x'\|_1=\tau
\quad\Longrightarrow\quad
d((s,x),(s,x'))\le C\tau\frac mN(\log m)^2.
\end{equation}
Moreover,
\begin{equation}\label{met:eq:size-diameter}
|V|\le\exp(Cm\log m),\qquad
\operatorname{diam}(V,d)\le Cm,
\end{equation}
and, for independent uniform $\theta,\eta\in T=(-1,1)^m$,
\begin{equation}\label{met:eq:average}
\E d(v_1(\theta),v_1(\eta))\ge cm.
\end{equation}
Here $N=m^6$, $\tau=m^{-2000}$, and all constants are absolute.
\end{proposition}

\begin{proof}
Let $P_v$ and $L=C_0\log m$ be as in \Cref{ker:angles}. Within a
chart set
\[
\delta((s,x),(s,x'))=L\frac mN\|x-x'\|_1,
\]
and put $\delta=\infty$ between different charts.
Equation \eqref{ker:global-angle} bounds every $P$-angle below by
$-r$. Equation \eqref{ker:local-angle} also bounds it below by
$-\delta(v,z)$ whenever $v,z$ are in the same chart, and by
$-\delta(y,z)$ whenever $y,z$ are in the same chart, after swapping
the two arguments of the inner product. Thus the hypothesis of
\Cref{lem:repair} holds. Use its vectors $q_v$ and set
\begin{equation}\label{met:eq:definition}
d(v,w)=|(P_v,q_v)-(P_w,q_w)|^2.
\end{equation}
This is a negative-type semimetric. The bound
$|q_v-q_w|^2\le8r$ and \eqref{ker:rounded-distance} prove
\eqref{met:eq:macro}.

Taking $y=v$ in \eqref{ker:local-angle} gives
$|P_v-P_z|^2\le\delta(v,z)$. If $\|x-x'\|_1=\tau$, then
\[
\delta= C_0\tau(m/N)\log m<r,
\qquad
\log(r/\delta)
=1997\log m-\log(C_0\log m)=O(\log m).
\]
Equation \eqref{rep:eq:local} now proves
\eqref{met:eq:adjacent}.

We next count the rounded labels. Since $|u_{s,i}|\le2$, a
hyperplane $\langle u_{s,i},\theta\rangle=k\tau$ meeting $U$ has
$|k\tau|<4\sqrt m$. Hence at most
\[
H\le N(1+8\sqrt m/\tau)
\]
threshold hyperplanes meet $U$ in a chart. An arrangement of $H$
hyperplanes in $\R^m$ has at most
$\sum_{j=0}^m\binom Hj$ full-dimensional regions (see~\cite{Zas75}).
This follows by
induction: an added hyperplane creates at most as many new regions
as the preceding hyperplanes cut out on it. The labels are constant
on the regions, and their intersections with the convex cube $U$
are connected. Therefore
\[
|V|\le S\sum_{j=0}^m\binom Hj\le\exp(Cm\log m),
\]
because $H$ is bounded by a fixed power of $m$. The diameter bound
in \eqref{met:eq:size-diameter} follows from
\eqref{met:eq:macro} and $\operatorname{diam}(U)=4\sqrt m$.

Finally, for independent uniform parameters on $T$,
\[
\E|\theta-\eta|^2=\frac{2m}{3},\qquad
|\theta-\eta|\le2\sqrt m.
\]
It follows that $\E|\theta-\eta|\ge\sqrt m/3$.
Equation \eqref{met:eq:macro} gives
$\E d(v_1(\theta),v_1(\eta))\ge c_*m/3-Cr$, proving
\eqref{met:eq:average} for sufficiently large $m$.
\end{proof}

The macroscopic comparison \eqref{met:eq:macro} does not make the
interface constraints redundant: its additive error $O(r)=O(m^{-8})$
is far larger than the interface bound $O(m^{-2005}(\log m)^2)$ in
\eqref{met:eq:adjacent}. The next section uses these much smaller
interface distances to constrain every contraction.
